\chapter{Finite-Volume Form Factors And Spectral Expansions}
\label{ch:integrable-finite-volume-form-factors}

\ConventionsLorentzian

Exact scattering data become QFT data only after they are connected to local
operator matrix elements and to correlation functions.  A spatial circle of
circumference \(L\) is the cleanest intermediate regulator in massive
two-dimensional integrable theories: rapidities become discrete, scattering
states become normalizable vectors, matrix elements have honest Hilbert-space
normalizations, and the infinite-volume spectral expansion is recovered as a
controlled limit.  This chapter develops that bridge for diagonal scattering
sectors.  Non-diagonal scattering requires replacing scalar phase functions by
eigenvalues of transfer matrices and is treated in the nested Bethe-ansatz
chapters.

\section{Finite-Volume Hilbert Space And Bethe--Yang Quantization}

Let the infinite-volume theory have stable particle species
\[
  a\in\mathcal I,
  \qquad
  p_a^\mu(\theta)=m_a(\cosh\theta,\sinh\theta),
\]
with diagonal two-body scattering eigenvalues
\[
  S_{ab}(\theta)=\exp(\ii\delta_{ab}(\theta)).
\]
The phase convention is fixed by unitarity and real analyticity on the real
rapidity axis:
\[
  S_{ab}(\theta)S_{ba}(-\theta)=1,
  \qquad
  \delta_{ab}(\theta)+\delta_{ba}(-\theta)\in2\pi\mathbb Z .
\]
We choose a continuous branch on the rapidity interval under consideration
and define
\[
  \varphi_{ab}(\theta)=\partial_\theta\delta_{ab}(\theta).
\]
The branch of \(\delta\) changes the Bethe integers, while
\(\varphi\) and the Gaudin density below are branch-independent.

For a multiparticle state with ordered rapidities
\[
  \theta_1,\ldots,\theta_n
\]
and species \(a_1,\ldots,a_n\), the Bethe--Yang equations are
\begin{equation}
  Q_j(\theta)
  =
  m_{a_j}L\sinh\theta_j
  +
  \sum_{k\ne j}
  \delta_{a_ja_k}(\theta_j-\theta_k)
  =
  2\pi I_j,
  \qquad j=1,\ldots,n .
  \label{eq:finite-volume-bethe-yang}
\end{equation}
The \(I_j\) are integers or half-integers depending on particle statistics,
twists, and branch choices.  The finite-volume energy and momentum at
Bethe--Yang order are
\[
  E_L(\theta)=\sum_{j=1}^n m_{a_j}\cosh\theta_j,
  \qquad
  P_L(\theta)=\sum_{j=1}^n m_{a_j}\sinh\theta_j .
\]
The neglected effects are exponentially small in \(m_\ast L\), where
\[
  m_\ast=\min_{a\in\mathcal I}m_a,
\]
provided no massless channel, threshold singularity, or nearby bound-state
pole invalidates the contour deformation used to wrap virtual particles
around the circle.

\begin{proposition}[Bethe--Yang phase from moving one particle around the circle]
\label{prop:bethe-yang-phase-around-circle}
Assume the \(n\)-particle wave function is described away from collision
hyperplanes by ordered plane waves, and that crossing a hyperplane
\(x_j=x_k\) multiplies the ordered amplitude by the diagonal scattering phase
\(S_{a_ja_k}(\theta_j-\theta_k)\).  Periodicity of particle \(j\) around the
circle gives \eqref{eq:finite-volume-bethe-yang}.
\end{proposition}

\begin{proof}
Move particle \(j\) once around the spatial circle while keeping the other
particles fixed.  The free propagation of this particle contributes
\[
  \exp(\ii p_jL),
  \qquad
  p_j=m_{a_j}\sinh\theta_j .
\]
During the circuit it crosses each other particle once.  In the diagonal
sector the ordered amplitude is multiplied by
\[
  \prod_{k\ne j}S_{a_ja_k}(\theta_j-\theta_k)
  =
  \exp\!\left(
    \ii\sum_{k\ne j}\delta_{a_ja_k}(\theta_j-\theta_k)
  \right).
\]
Single-valuedness of the wave function demands that the total phase be one.
Taking the logarithm on the chosen branch gives
\[
  m_{a_j}L\sinh\theta_j
  +
  \sum_{k\ne j}\delta_{a_ja_k}(\theta_j-\theta_k)
  =
  2\pi I_j .
\]
\end{proof}

\section{Gaudin Density And Normalization}

The density of Bethe--Yang states is the Jacobian of the map
\[
  \theta\longmapsto Q(\theta).
\]
Define the Gaudin matrix
\begin{equation}
  G_{ij}(\theta)
  =
  \frac{\partial Q_i}{\partial\theta_j},
  \qquad
  \rho_n(\theta)=\det G(\theta).
  \label{eq:gaudin-matrix-density}
\end{equation}
Explicitly,
\[
  G_{ii}
  =
  m_{a_i}L\cosh\theta_i
  +
  \sum_{k\ne i}\varphi_{a_ia_k}(\theta_i-\theta_k),
  \qquad
  G_{ij}
  =
  -\varphi_{a_ia_j}(\theta_i-\theta_j)
  \quad (i\ne j).
\]
For one particle,
\[
  \rho_1(\theta)=m_aL\cosh\theta .
\]
For two particles, using
\(\varphi_{a_2a_1}(\theta_2-\theta_1)=
\varphi_{a_1a_2}(\theta_1-\theta_2)\), one obtains
\begin{equation}
  \rho_2(\theta_1,\theta_2)
  =
  m_{a_1}m_{a_2}L^2\cosh\theta_1\cosh\theta_2
  +
  \varphi_{a_1a_2}(\theta_{12})
  L\bigl(
    m_{a_1}\cosh\theta_1+m_{a_2}\cosh\theta_2
  \bigr),
  \label{eq:two-particle-gaudin-density}
\end{equation}
where \(\theta_{12}=\theta_1-\theta_2\).

\begin{proposition}[State-counting measure]
\label{prop:gaudin-state-counting-measure}
Let \(U\subset\mathbb R^n\) be a rapidity region in which the Bethe--Yang map
\(Q:U\to\mathbb R^n\) is one-to-one and \(\rho_n>0\).  For every smooth
compactly supported test function \(f\) in \(U\),
\[
  \sum_{I:\,\theta(I)\in U}f(\theta(I))
  =
  \int_U
  \frac{\dd^n\theta}{(2\pi)^n}\,
  \rho_n(\theta)f(\theta)
  +
  O(L^{n-1})
\]
as \(L\to\infty\), with the boundary correction depending on \(f\) and \(U\).
\end{proposition}

\begin{proof}
The Bethe integers form the lattice \((2\pi)\mathbb Z^n\) or a translate of
it in \(Q\)-space.  Since \(Q\) is one-to-one on \(U\), summing over Bethe
solutions in \(U\) is summing the function
\[
  f(\theta(Q))
\]
over lattice points in \(Q(U)\).  The leading Euler--Maclaurin term is the
Lebesgue integral in \(Q\)-space with lattice cell volume \((2\pi)^n\):
\[
  \sum_I f(\theta(I))
  =
  \int_{Q(U)}
  \frac{\dd^nQ}{(2\pi)^n}\,
  f(\theta(Q))
  +
  O(\operatorname{area}\partial Q(U)).
\]
Changing variables from \(Q\) to \(\theta\) gives
\[
  \dd^nQ=\det\left(\frac{\partial Q_i}{\partial\theta_j}\right)\dd^n\theta
  =
  \rho_n(\theta)\dd^n\theta .
\]
The boundary term is \(O(L^{n-1})\) because \(Q(U)\) has linear size
\(O(L)\) in each direction for fixed \(U\).
\end{proof}

The same determinant normalizes finite-volume states.  If
\[
  \ket{\theta_1,a_1;\ldots;\theta_n,a_n}_{\infty}
\]
is an infinite-volume rapidity state normalized by delta functions, then the
corresponding unit finite-volume vector has leading normalization
\[
  \ket{I_1,\ldots,I_n}_L
  =
  \rho_n(\theta)^{-1/2}
  \ket{\theta_1,a_1;\ldots;\theta_n,a_n}_{\infty}
  +
  O(\ee^{-\mu L})
\]
in matrix elements of local operators.  The statement compares normalized
matrix elements through the finite-volume regulator; the infinite-volume
delta-normalized state and the finite-volume unit vector live in different
Hilbert-space normalizations.

\section{Off-Diagonal Matrix Elements}

Let \(\widehat{\mathcal O}\) be a local scalar operator.  The infinite-volume
crossed form factor for an outgoing \(m\)-particle state and incoming
\(n\)-particle state is written
\[
  F^{\mathcal O}_{b_1\ldots b_m,a_1\ldots a_n}
  (\theta'_1+\ii\pi,\ldots,\theta'_m+\ii\pi,
    \theta_1,\ldots,\theta_n),
\]
with the branch specified by the rapidity ordering and by the crossing
prescription of Chapter~\ref{ch:integrable-form-factor-bootstrap-local-operators}.

\begin{proposition}[Off-diagonal finite-volume form-factor formula]
\label{prop:off-diagonal-finite-volume-form-factor}
Assume the incoming and outgoing Bethe--Yang rapidity sets have no coincident
rapidities after crossing, and assume that the crossed form factor is regular
at the corresponding point.  Then
\begin{equation}
  {}_L\!\langle J_1,\ldots,J_m|
  \widehat{\mathcal O}(0)
  |I_1,\ldots,I_n\rangle_L
  =
  \frac{
    F^{\mathcal O}_{b_1\ldots b_m,a_1\ldots a_n}
    (\theta'_1+\ii\pi,\ldots,\theta'_m+\ii\pi,
      \theta_1,\ldots,\theta_n)
  }{
    \sqrt{\rho_m(\theta')\rho_n(\theta)}
  }
  +
  O(\ee^{-\mu L}).
  \label{eq:off-diagonal-finite-volume-form-factor}
\end{equation}
\end{proposition}

\begin{proof}
Smear the finite-volume Bethe states by packets whose rapidity widths are
large compared with the level spacing \(L^{-1}\) and small compared with the
distance to any kinematic or bound-state pole.  In that window the
finite-volume packet is approximated by the infinite-volume packet with the
state-counting measure of Proposition~\ref{prop:gaudin-state-counting-measure}.
The unit normalization of the finite-volume packet therefore supplies one
factor \(\rho_n^{-1/2}\) for the incoming state and one factor
\(\rho_m^{-1/2}\) for the outgoing state.  The local operator matrix element
between the corresponding infinite-volume packets is the crossed form factor
by the crossing prescription.  Since the form factor is regular at the
chosen point, shrinking the packets to the Bethe rapidities gives
\eqref{eq:off-diagonal-finite-volume-form-factor}.  Virtual particles
wrapping around the circle give the exponentially small correction.
\end{proof}

\begin{corollary}[Cancellation in the infinite-volume two-point sum]
\label{cor:finite-volume-sum-integral-cancellation}
For vacuum-to-\(n\)-particle matrix elements of a scalar operator,
\[
  \sum_I
  \left|
  {}_L\!\langle\Omega_L|\widehat{\mathcal O}(0)|I\rangle_L
  \right|^2 f(\theta(I))
  \longrightarrow
  \int
  \frac{\dd^n\theta}{(2\pi)^n}
  |F_n^{\mathcal O}(\theta)|^2f(\theta)
\]
on rapidity regions with no singularities, after the usual \(1/n!\) quotient
for identical unordered particles.
\end{corollary}

\begin{proof}
By Proposition~\ref{prop:off-diagonal-finite-volume-form-factor}, the squared
matrix element contributes
\[
  |F_n^{\mathcal O}(\theta(I))|^2/\rho_n(\theta(I))
\]
up to exponentially small corrections.  Summing over Bethe integers and
using Proposition~\ref{prop:gaudin-state-counting-measure} supplies the
measure
\[
  \rho_n(\theta)\dd^n\theta/(2\pi)^n .
\]
The Gaudin determinant cancels.  If the rapidity ordering was used to label
states, the ordered chamber gives one copy; if integration is over all
\(\mathbb R^n\) for identical particles, the \(n!\) permutations must be
divided out.
\end{proof}

\section{Diagonal Matrix Elements}

Diagonal matrix elements require subtracting kinematic poles.  Let
\[
  F^{\mathcal O}_{2n,c}(\theta_1,\ldots,\theta_n)
\]
denote the connected diagonal form factor obtained as the finite part of
\[
  F^{\mathcal O}_{\bar a_n\ldots\bar a_1a_1\ldots a_n}
  (\theta_n+\ii\pi+\epsilon_n,\ldots,\theta_1+\ii\pi+\epsilon_1,
  \theta_1,\ldots,\theta_n)
\]
after removing all terms singular in the \(\epsilon_j\) as
\(\epsilon_j\to0\).  The order of the crossed rapidities is part of the
definition.  In diagonal theories the connected finite part is symmetric in
the rapidities after the same branch convention is used throughout.

\begin{definition}[Principal minors of the Gaudin matrix]
\label{def:gaudin-principal-minors}
For \(A\subset\{1,\ldots,n\}\), let \(G_A\) be the principal submatrix of
the Gaudin matrix with rows and columns in \(A\), and set
\[
  \rho_A(\theta)=\det G_A,
  \qquad
  \rho_\emptyset=1 .
\]
When the complement of \(A\) is denoted \(\bar A\), the symbol
\(\rho_{\bar A}\) means the determinant of the principal minor on that
complement, not the Gaudin density of an isolated smaller-volume theory.
\end{definition}

\begin{proposition}[Diagonal finite-volume matrix element]
\label{prop:diagonal-finite-volume-form-factor}
Under the same massive diagonal assumptions and away from accidental
degeneracies,
\begin{equation}
  {}_L\!\langle I_1,\ldots,I_n|
  \widehat{\mathcal O}(0)
  |I_1,\ldots,I_n\rangle_L
  =
  \frac{1}{\rho_n(\theta)}
  \sum_{A\subset\{1,\ldots,n\}}
  F^{\mathcal O}_{2|A|,c}(\theta_A)\,
  \rho_{\bar A}(\theta)
  +
  O(\ee^{-\mu L}).
  \label{eq:diagonal-finite-volume-form-factor}
\end{equation}
Here \(F^{\mathcal O}_{0,c}\) is the infinite-volume vacuum expectation value
\(\langle\mathcal O\rangle\).
\end{proposition}

\begin{proof}
Start from the off-diagonal formula with outgoing rapidities
\(\theta'_j\) distinct from the incoming rapidities, then approach
\(\theta'_j=\theta_j\).  The crossed form factor develops kinematic poles.
Each singular contraction identifies one crossed particle with one incoming
particle and produces the finite-volume delta contribution associated with
the corresponding Bethe quantum number.  In the finite-volume regularization
these delta contributions are replaced by derivatives of the quantization
conditions, namely entries and minors of the Gaudin matrix.

After all possible disconnected contractions are grouped by the subset
\(\bar A\) of particles that contract through the finite-volume
normalization, the remaining particles \(A\) contribute the connected
diagonal finite part \(F^{\mathcal O}_{2|A|,c}(\theta_A)\), while the
contracted particles contribute the principal minor \(\rho_{\bar A}\).  The
overall normalization of the bra and ket contributes \(1/\rho_n\).  Summing
over all subsets gives \eqref{eq:diagonal-finite-volume-form-factor}.
The only analytic input beyond the off-diagonal case is the kinematic-pole
residue equation of Chapter~\ref{ch:integrable-form-factor-bootstrap-local-operators},
which ensures that the disconnected pieces are exactly the contractions
enumerated by the subsets.
\end{proof}

For one particle this gives
\[
  {}_L\!\langle I|\mathcal O|I\rangle_L
  =
  \langle\mathcal O\rangle
  +
  \frac{F^{\mathcal O}_{2,c}(\theta)}{\rho_1(\theta)}
  +
  O(\ee^{-\mu L}).
\]
For two particles,
\[
\begin{aligned}
  {}_L\!\langle I_1,I_2|\mathcal O|I_1,I_2\rangle_L
  &=
  \frac{1}{\rho_{12}}
  \Bigl[
    \langle\mathcal O\rangle\rho_{12}
    +F_{2,c}^{\mathcal O}(\theta_1)\rho_{\{2\}}
    +F_{2,c}^{\mathcal O}(\theta_2)\rho_{\{1\}}  \\
  &\hspace{3.4cm}
    +F_{4,c}^{\mathcal O}(\theta_1,\theta_2)
  \Bigr]
  +
  O(\ee^{-\mu L}).
\end{aligned}
\]
This explicit formula is often the first place where the distinction between
connected diagonal form factors and symmetric finite parts matters.

\section{Zero-Temperature Spectral Expansion}

For a finite-volume two-point function
\[
  C_L(x,t)=
  \langle\Omega_L|
  \widehat{\mathcal O}(x,t)\widehat{\mathcal O}(0,0)
  |\Omega_L\rangle ,
\]
insert the finite-volume identity operator:
\[
  \mathbf 1_L=\sum_\alpha|\alpha,L\rangle\langle\alpha,L|.
\]
Translation covariance gives
\[
  C_L(x,t)
  =
  \sum_\alpha
  \bigl|
    \langle\Omega_L|\widehat{\mathcal O}(0)|\alpha,L\rangle
  \bigr|^2
  \exp(-\ii E_\alpha t+\ii P_\alpha x).
\]
Using Corollary~\ref{cor:finite-volume-sum-integral-cancellation} in each
particle-number sector gives the infinite-volume Wightman expansion
\begin{equation}
  C(x,t)
  =
  \sum_{n=0}^{\infty}
  \frac{1}{n!}
  \sum_{a_1,\ldots,a_n}
  \int_{\mathbb R^n}
  \prod_{j=1}^n\frac{\dd\theta_j}{2\pi}\,
  \bigl|
    F^{\mathcal O}_{a_1\ldots a_n}(\theta_1,\ldots,\theta_n)
  \bigr|^2
  e^{-\ii\sum_j m_{a_j}t\cosh\theta_j
      +\ii\sum_j m_{a_j}x\sinh\theta_j}.
  \label{eq:zero-temperature-form-factor-spectral-expansion}
\end{equation}
The Euclidean Schwinger function at separation
\[
  r=\sqrt{\tau^2+x^2}>0
\]
is obtained by analytic continuation \(t=-\ii\tau\).  Lorentz invariance
allows the exponent to be written, in a frame where the separation is purely
Euclidean time, as
\[
  \exp\left[-r\sum_{j=1}^n m_{a_j}\cosh\theta_j\right].
\]

\begin{proposition}[A sufficient convergence criterion]
\label{prop:form-factor-spectral-convergence-criterion}
Fix \(r>0\).  Suppose there are constants \(A,B,c\ge0\) with \(c<m_\ast r\)
such that, for every \(n\),
\[
  \frac{1}{n!}
  \sum_{a_1,\ldots,a_n}
  \int_{\mathbb R^n}
  \prod_{j=1}^n\frac{\dd\theta_j}{2\pi}\,
  |F^{\mathcal O}_{a_1\ldots a_n}(\theta)|^2
  e^{-r\sum_jm_{a_j}\cosh\theta_j}
  \le
  \frac{A B^n}{n!}e^{cn}.
\]
Then the Euclidean spectral series converges absolutely.
\end{proposition}

\begin{proof}
The right hand side is summable in \(n\):
\[
  \sum_{n=0}^\infty \frac{A(B e^c)^n}{n!}
  =
  A\exp(Be^c).
\]
The hypothesis bounds the absolute value of the \(n\)-particle contribution
by a summable sequence, so the series converges absolutely by comparison.
The displayed condition is deliberately a sufficient condition; sharp
convergence in specific integrable models is a model-dependent analytic
problem involving the large-rapidity and large-particle-number behavior of
the form factors.
\end{proof}

\section{Example: Free Majorana Energy Density}

In the massive Ising field theory without magnetic perturbation, the free
Majorana scattering phase is
\[
  S=-1.
\]
For the energy-density operator \(\varepsilon\), the only nonzero
vacuum-to-particle form factor in the even sector needed for the first
spectral term is the two-particle form factor
\[
  F_2^\varepsilon(\theta_1,\theta_2)
  =
  \kappa\sinh\left(\frac{\theta_1-\theta_2}{2}\right),
\]
with normalization \(\kappa\) fixed by the convention for
\(\varepsilon\).  The two-particle contribution to the Euclidean two-point
function is
\begin{equation}
  C_2^\varepsilon(r)
  =
  \frac12
  \int_{\mathbb R^2}
  \frac{\dd\theta_1\dd\theta_2}{(2\pi)^2}
  \left|
    \kappa\sinh\frac{\theta_1-\theta_2}{2}
  \right|^2
  e^{-mr(\cosh\theta_1+\cosh\theta_2)} .
  \label{eq:ising-energy-two-particle-spectral-term}
\end{equation}
Introduce center and relative rapidities
\[
  \Theta=\frac{\theta_1+\theta_2}{2},
  \qquad
  \alpha=\theta_1-\theta_2.
\]
Then
\[
  \cosh\theta_1+\cosh\theta_2
  =
  2\cosh\Theta\cosh\frac{\alpha}{2},
  \qquad
  \dd\theta_1\dd\theta_2=\dd\Theta\,\dd\alpha.
\]
Using
\[
  \int_{-\infty}^{\infty}
  e^{-z\cosh\Theta}\,\dd\Theta
  =
  2K_0(z),
\]
\eqref{eq:ising-energy-two-particle-spectral-term} becomes
\[
  C_2^\varepsilon(r)
  =
  \frac{|\kappa|^2}{(2\pi)^2}
  \int_{-\infty}^{\infty}
  \dd\alpha\,
  \sinh^2\frac{\alpha}{2}\,
  K_0\!\left(2mr\cosh\frac{\alpha}{2}\right).
\]
Since the integrand is even,
\[
  C_2^\varepsilon(r)
  =
  \frac{2|\kappa|^2}{(2\pi)^2}
  \int_0^\infty
  \dd\alpha\,
  \sinh^2\frac{\alpha}{2}\,
  K_0\!\left(2mr\cosh\frac{\alpha}{2}\right).
\]
This example displays the finite-volume logic in its simplest form: the
finite-volume matrix element contains the inverse Gaudin determinant, the sum
over Bethe integers contains the Gaudin determinant, and the infinite-volume
answer contains only the invariant rapidity measure and the form factor.

\section{Thermal Traces And Connected Diagonal Series}

At finite temperature, the Gibbs trace is organized by the thermodynamic
limit at fixed rapidity density.  The pseudoenergy
\(\varepsilon_a(\theta)\) is determined by the TBA equations in the same
scattering convention used to define the form factors.
For fermionic Bethe quantization conventions, the occupation factor is
\[
  n_a(\theta)=\frac{1}{1+e^{\varepsilon_a(\theta)}}.
\]
The Leclair--Mussardo type one-point series has the form
\begin{equation}
  \langle\mathcal O\rangle_R
  =
  \sum_{n=0}^{\infty}\frac{1}{n!}
  \sum_{a_1,\ldots,a_n}
  \int\prod_{j=1}^n\frac{\dd\theta_j}{2\pi}\,
  \prod_{j=1}^n n_{a_j}(\theta_j)\,
  F^{\mathcal O}_{2n,c}(\theta_1,\ldots,\theta_n),
  \label{eq:lm-connected-diagonal-series}
\end{equation}
where \(R\) is inverse temperature.  Other quantization conventions move
statistical signs between the occupation factor and the connected diagonal
form factor.  The formula is meaningful only with a fixed prescription for
the TBA kernel, the connected diagonal finite part, and the convergence
domain.

\section{Reconstruction Boundary}

The finite-volume form-factor formalism supplies a computable chain
\[
  S\text{-matrix}
  \longrightarrow
  F_n^{\mathcal O}
  \longrightarrow
  {}_L\!\langle\beta|\mathcal O|\alpha\rangle_L
  \longrightarrow
  \text{correlation functions}.
\]
It does not by itself prove that the resulting distributions define a local
QFT.  A construction theorem must prove positivity of the reconstructed
correlation functions, covariance, locality or wedge-locality with a local
subalgebra, clustering, and compatibility with the scattering theory used at
the beginning.  The finite-volume regulator is valuable precisely because
each pre-limit step has a declared Hilbert space and normalizable vectors.

\begin{openproblem}[Form-factor construction]
\label{op:form-factor-construction}
Construct an interacting factorizing QFT from bootstrap scattering and
form-factor data by proving convergence of the spectral expansions, locality
of reconstructed fields or algebras, and agreement of finite-volume matrix
elements with the Hilbert-space scattering limit.
\end{openproblem}
